% !TeX program = xelatex
% SPDX-License-Identifier: MIT
% Copyright (c) 2026 Ziyu Peng and 刘毅涵
% Author: https://pzy54154631-hub.github.io/homepage/
% Author: https://pzy54154631-hub.github.io/liuyihan/
% 自设教学现金流；不对应任何真实投资产品。
\documentclass[UTF8,a4paper,fontset=fandol]{ctexart}
\usepackage[margin=2.5cm]{geometry}
\usepackage{amsmath,amssymb,booktabs,siunitx}
\sisetup{group-separator={,},group-minimum-digits=4}
\usepackage[hidelinks]{hyperref}
\title{金融公式：让现金流与时间对齐}
\author{\href{https://pzy54154631-hub.github.io/homepage/}{\textit{Ziyu Peng}}, University of Colorado Boulder
  \\[0.35em]
  \href{https://pzy54154631-hub.github.io/liuyihan/}{刘毅涵}，暨南大学经济学院}
\date{}
\begin{document}
\maketitle
\section{收益率与复利}
无分红、无中途现金流、价格均为正时，简单收益率与对数收益率为
\[
 R_t=\frac{P_t-P_{t-1}}{P_{t-1}},\qquad
 r_t=\ln\!\left(\frac{P_t}{P_{t-1}}\right)
     =\ln(1+R_t).
\]
价格从 100 变为 110 时，$R_t=10\%$，$r_t\approx0.09531$。
若每期有效利率恒为 $i$，且收益在相同利率下再投资，则
\[
 \mathrm{FV}_n=\mathrm{PV}_0(1+i)^n.
\]
例如本金 1000、年有效利率 $5\%$、期限 3 年，有
$\mathrm{FV}_3=1000(1.05)^3=1157.625$，显示到分为 1157.63。

\section{净现值：明确每笔现金流的时点}
设期初 $t=0$ 支出 1000，第 1、2、3 年末分别收到 400、450、500，
年有效折现率取 $8\%$。金额单位统一为元，参数均为教学设定。
\begin{align*}
 \mathrm{NPV}
 &=\sum_{t=0}^{3}\frac{\mathrm{CF}_t}{(1+i)^t}\\
 &=-1000+\frac{400}{1.08}
   +\frac{450}{1.08^2}+\frac{500}{1.08^3}\\
 &\approx153.09.
\end{align*}
\begin{center}
\begin{tabular}{c S[table-format=-4.0] S[table-format=-4.2]}
\toprule
{时点 $t$（年）} & {现金流（元）} & {现值（元）}\\
\midrule
0 & -1000 & -1000.00\\
1 & 400 & 370.37\\
2 & 450 & 385.80\\
3 & 500 & 396.92\\
\midrule
\multicolumn{2}{r}{NPV（由未舍入值求和）} & 153.09\\
\bottomrule
\end{tabular}
\end{center}
表中单行显示值已舍入；计算 NPV 时使用未舍入值，最后统一显示到分。
$t=0$ 表示期初；$t=1,2,3$ 分别表示第 1、2、3 年末。

\section{固定息票债券的课堂估值}
设面值 $F=1000$、年息票率 $c=4\%$，每年末付息一次，
3 年后还本；估值时点刚付完息，年有效收益率假设为 $y=5\%$。
忽略违约、税费与嵌入期权，则
\[
 B_0=\sum_{t=1}^{3}\frac{cF}{(1+y)^t}
     +\frac{F}{(1+y)^3}
     \approx972.77.
\]
最后一期收到的总额为 $40+1000=1040$，不能只计息票。
本例仅用于演示排版与贴现计算。
\end{document}
